\chapter{Implementación de las simulaciones}\label{ap:simulacion}

Todas las simulaciones del capítulo~\ref{chap:desempeno} siguen el mismo esquema de tiempo discreto, que reproduce un controlador digital con retenedor de orden cero. En cada periodo de muestreo $T=1$~ms se ejecutan los pasos siguientes.

\begin{enumerate}
  \item \textbf{Planta.} Se integra el modelo no lineal durante un periodo con el voltaje $u$ constante, con un integrador numérico de paso variable para ecuaciones diferenciales ordinarias. Con atascamiento-deslizamiento, la aceleración se calcula con el modelo de Karnopp: si $|v|>DV$, el imán desliza con fricción viscosa $m c_1 v$; si $|v|\le DV$ y la fuerza neta $F_e=u/[b(a-x)^4]-mg$ cumple $|F_e|\le F_s$, permanece atascado; en otro caso se despega con aceleración $(F_e - F_s\,\operatorname{sign}F_e)/m$.
  \item \textbf{Sujeción de atascamiento.} Al final del periodo, si $|v|<DV$ y $|F_e|\le F_s$, la velocidad se fija en cero. Esto evita que el integrador numérico deje una velocidad residual dentro de la banda de Karnopp.
  \item \textbf{Controlador.} Se calcula el error $e=r-x$ y se integra el modelo en espacio de estados del controlador durante el mismo periodo con $e$ constante.
  \item \textbf{Saturación.} La salida del controlador se acota al intervalo $[0;\,3.5]$~V. Las fracciones de saturación que se reportan se miden sobre el voltaje aplicado, después del paso siguiente.
  \item \textbf{Zona muerta.} Si el imán está en movimiento y el nuevo voltaje difiere del último aplicado en menos de $+0.025$~V o $-0.02$~V, se conserva el último voltaje aplicado.
\end{enumerate}

El controlador se inicializa en su estado estacionario: se resuelve $(A_d - I)\,x_c = 0$ junto con $C x_c = u_e(x_0)$, donde $A_d$ es la matriz del controlador discretizado, de modo que en $t=0$ entrega exactamente el voltaje de equilibrio de la posición inicial y el levitante parte del reposo.

Los barridos paramétricos (rango viable del actuador y familias de controladores) se ejecutaron con una implementación equivalente en la GPU, en precisión simple, con el controlador discretizado de forma exacta y la planta integrada por Runge-Kutta de cuarto orden con diez subpasos por periodo. Las 512 trayectorias de las familias se repitieron en doble precisión con diez y con veinte subpasos. La diferencia mediana del ciclo límite respecto a la GPU es inferior al 1\,\%, y la mayor es del 16\,\% en un solo controlador, que la integración de paso variable con tolerancias estrictas confirma (sección~\ref{sec:barrido_familias}). Los resultados individuales del barrido tienen, por tanto, esa incertidumbre numérica; los de las demás secciones provienen de esa misma implementación de paso variable.

La tabla~\ref{tab:scripts} relaciona cada procedimiento de cálculo con los resultados de la tesis que genera.

\begin{table}[ht]
\centering
\caption{Procedimientos de cálculo que generan los resultados de la tesis.}
\label{tab:scripts}
\small
\begin{tabular}{|p{0.42\textwidth}|p{0.48\textwidth}|}
\hline
Procedimiento & Resultados \\ \hline
Evidencia de lazo abierto & Lazo abierto, modelo lineal con AD, configuración estable \\ \hline
Diseño del PII & Diagramas de Bode, Nyquist y respuesta al escalón \\ \hline
Regulación sin atascamiento-deslizamiento & Modelo no lineal sin atascamiento-deslizamiento \\ \hline
Regulación & Respuesta ante referencias escalón \\ \hline
Seguimiento & Seguimiento senoidal y trapezoidal \\ \hline
Comparación PI frente a PII & Comparación directa de ambos controladores en los mismos escenarios \\ \hline
Familias de controladores & Familias de controladores PI y PII \\ \hline
Exportación de familias & Criterio de aceptación y datos de la figura de familias \\ \hline
Análisis del ciclo límite & Mecanismo y amplitud del ciclo límite \\ \hline
Barrido del rango viable & Rango viable con $u\le3.5$~V \\ \hline
\end{tabular}
\end{table}
